\subsection{\texorpdfstring{From CB\textsuperscript{2} to
BARCS}{From CB2 to BARCS}}

CB\textsuperscript{2} adapts the beta-binomial weighted test of
\citet{baggerly2003differential} to CRISPR pooled screens
\citep{jeong2019beta}.  Its central insight is worth preserving: sequencing
depth measures a guide proportion precisely, but it does not turn reads into
independent biological replicates.  The original method applies this insight
to a two-group contrast.  BARCS retains the same
depth-aware, extra-binomial variance principle and changes the mean model from
two group means to a design matrix.  ``Retains'' refers to the stochastic
hierarchy and the sample-level information ceiling; it does not mean that the
default logit fit reproduces the original finite-sample statistic.  The
legacy representation and asymptotic relationship are established below.

\begin{table}[ht]
\centering
\caption{What changes---and what does not---from the original
CB\textsuperscript{2} method to BARCS.}
\label{tab:cb2difference}
\small
\setlength{\tabcolsep}{5pt}
\begin{tabular}{p{0.20\textwidth}p{0.32\textwidth}p{0.38\textwidth}}
\toprule
Feature & Original CB\textsuperscript{2} & BARCS\\
\midrule
Primary question & Difference between two groups &
Trend, adjusted association, interaction, or contrast\\
Sample design & Reference and comparison labels &
Arbitrary full-rank R model matrix\\
Mean model & Two unrelated group proportions &
$\logit(\mu_i)=\mathbf{x}_i^\top\boldsymbol{\beta}$\\
Effect & Difference or fold change between groups &
Conditional log-odds coefficient or linear contrast\\
Dispersion & Separate beta-binomial weighting within each group &
One guide-specific $\rho$ shared across the fitted design\\
Estimator & Weighted group proportions &
Feasible logistic IRLS\\
Variance principle & \multicolumn{2}{c}{Binomial sequencing plus
between-library beta-binomial heterogeneity}\\
Reference df & Group-variance df formula &
Residual $m-q$\\
Adjustment & Pairwise stratification or preprocessing &
Batch, donor, dose, time, interactions, splines, and other covariates\\
Implementation & Existing CB\textsuperscript{2} workflow &
Additive \texttt{bbreg()}, \texttt{bb\_contrast()}, and
\texttt{bb\_screen()} APIs with RcppArmadillo kernels\\
\bottomrule
\end{tabular}
\end{table}

The regression bridge was described by \citet{baggerly2004overdispersed} for
sequencing count data with multiple groups and continuous covariates.  The
method combines logistic regression with extra-binomial variation and retains
a coefficient-based $t$ test.  Here we state that construction in CRISPR-screen
notation and use the beta-binomial, library-size-dependent variance of
\citet{williams1982extrabinomial}.  An identity-link regression of raw
proportions would allow fitted values outside $(0,1)$, while an unweighted
correlation would ignore unequal library sizes.  The logit and beta-binomial
weights solve those two problems directly.

\subsection{Model}

For guide $g$ in sequenced sample $i$, let $K_{gi}$ be its read count, let
$N_i$ be the full mapped-guide total for that sample, and let
$Y_{gi}=K_{gi}/N_i$ be the observed guide proportion.  Sample-level
quantities such as dose, time, treatment, donor, and batch are predictors even
when they are described biologically as experimental ``outcomes.''

The model is fitted separately for every guide, so the guide index is
suppressed below.  Conditional on a latent sample-specific abundance $P_i$,
\begin{align}
  K_i \mid P_i &\sim \operatorname{Binomial}(N_i,P_i),\\
  P_i &\sim \operatorname{Beta}(\mu_i\kappa,(1-\mu_i)\kappa).
\end{align}
The mean is linked to a row vector of sample covariates $\mathbf{x}_i^\top$ by
\begin{equation}
  \logit(\mu_i)
  = \log\left(\frac{\mu_i}{1-\mu_i}\right)
  = \mathbf{x}_i^\top\boldsymbol{\beta}.
  \label{eq:mean}
\end{equation}
The design can contain an intercept, a continuous phenotype, batch indicators,
interactions, spline terms, and other prespecified adjustment variables.

It is useful to express the beta precision $\kappa$ as an intraclass
correlation
\begin{equation}
  \rho = \frac{1}{\kappa+1}, \qquad 0\leq\rho<1.
\end{equation}
Marginally,
\begin{align}
  \E(K_i) &= N_i\mu_i,\\
  \Var(K_i) &=
  N_i\mu_i(1-\mu_i)\{1+(N_i-1)\rho\}, \label{eq:countvar}\\
  \Var(Y_i) &=
  \mu_i(1-\mu_i)
  \left\{\frac{1-\rho}{N_i}+\rho\right\}. \label{eq:propvar}
\end{align}
Equation~\eqref{eq:propvar} cleanly separates sequencing uncertainty, which
decreases with $N_i$, from between-sample heterogeneity, which does not vanish
when sequencing depth becomes arbitrarily large.

\subsubsection{Interpretation of the continuous coefficient}

Suppose the model contains a continuous phenotype $d_i$ with coefficient
$\beta_d$.  Then $\exp(\beta_d)$ is the multiplicative change in the odds of
observing that guide for a one-unit increase in $d_i$, conditional on the other
covariates.  Standardizing $d_i$ before fitting makes $\exp(\beta_d)$ the odds
ratio per sample standard deviation and makes guide effects directly
comparable.  For rare guides, odds and proportions are similar, so $\beta_d$ is
also approximately a log abundance ratio per unit of $d_i$.

\subsection{Estimation by feasible IRLS}

Let $X$ be the $m\times q$ full-rank design matrix.  Given current values of
$\boldsymbol{\beta}$ and $\rho$, define
\begin{align}
  \eta_i &= \mathbf{x}_i^\top\boldsymbol{\beta}, &
  \mu_i &= \frac{\exp(\eta_i)}{1+\exp(\eta_i)},\\
  z_i &= \eta_i+\frac{Y_i-\mu_i}{\mu_i(1-\mu_i)}, &
  w_i &= \frac{N_i\mu_i(1-\mu_i)}
  {1+(N_i-1)\rho}. \label{eq:irls}
\end{align}
The coefficient update is the weighted least-squares step
\begin{equation}
  \widehat{\boldsymbol{\beta}}_{\mathrm{new}}
  =(X^\top W X)^{-1}X^\top W\mathbf{z},
  \qquad W=\operatorname{diag}(w_1,\ldots,w_m).
  \label{eq:betaupdate}
\end{equation}
This is weighted least squares on the logistic working response, not on the raw
proportions.

For a fixed fitted mean, $\rho$ is estimated from the Pearson equation
\begin{equation}
  \sum_{i=1}^{m}
  \frac{(K_i-N_i\widehat{\mu}_i)^2}
  {N_i\widehat{\mu}_i(1-\widehat{\mu}_i)
  \{1+(N_i-1)\widehat{\rho}\}}
  =m-q. \label{eq:rho}
\end{equation}
If the left side at $\rho=0$ is no larger than $m-q$, the estimate is truncated
at zero rather than claiming negative overdispersion.  Equations
\eqref{eq:irls}--\eqref{eq:rho} are alternated to convergence.

The weight has an important limiting behavior:
\begin{equation}
  \lim_{N_i\rightarrow\infty} w_i
  = \frac{\mu_i(1-\mu_i)}{\rho}
  \quad\text{when }\rho>0.
\end{equation}
Therefore a deeply sequenced sample cannot acquire unlimited leverage.  Once
sequencing uncertainty is small, biological heterogeneity sets the information
ceiling.

\subsection{T-based inference}

At convergence, the model-based covariance estimator is
\begin{equation}
  \widehat{\operatorname{Cov}}
  (\widehat{\boldsymbol{\beta}})
  =(X^\top\widehat{W}X)^{-1}.
  \label{eq:covariance}
\end{equation}
The Pearson equation scales the residual variation to approximately one.  If
the dispersion estimate reaches its numerical upper boundary, the
implementation additionally multiplies Equation~\eqref{eq:covariance} by the
remaining Pearson scale.

Equation~\eqref{eq:covariance} is a plug-in, model-based covariance: it
conditions on the estimated $\widehat{\rho}$ as though it were fixed.  It does
not propagate uncertainty from estimating a separate dispersion for every
guide.  Referring the coefficient to $t_{m-q}$ gives a heavier-tailed
small-sample reference than a normal approximation, but it is not a formal
correction for dispersion-estimation uncertainty.  With few independent
libraries, calibration must therefore be checked by simulation, phenotype
permutation, or prespecified negative controls.

For coefficient $\beta_j$, the statistic
\begin{equation}
  t_j =
  \frac{\widehat{\beta}_j-\beta_{j,0}}
  {\sqrt{
  \widehat{\operatorname{Cov}}
  (\widehat{\boldsymbol{\beta}})_{jj}}}
  \label{eq:tstat}
\end{equation}
is compared with a Student $t_{m-q}$ distribution.  The degrees of freedom are
driven by the number of independent sequenced samples, not by the total read
count.  This is the crucial reason that a library with millions of reads does
not create artificial certainty.

More generally, for a prespecified contrast vector $\mathbf{c}$,
\begin{equation}
  t_{\mathbf{c}} =
  \frac{\mathbf{c}^\top\widehat{\boldsymbol{\beta}}-\delta_0}
  {\sqrt{\mathbf{c}^\top
  \widehat{\operatorname{Cov}}(\widehat{\boldsymbol{\beta}})
  \mathbf{c}}}
\end{equation}
uses the same $m-q$ degrees of freedom.  This covers pairwise group contrasts,
average slopes in interaction models, and other one-degree-of-freedom
hypotheses.  Guide-wise two-sided $p$-values are adjusted using the
Benjamini--Hochberg procedure \citep{benjamini1995controlling}.

\subsubsection{Legacy representation and local-equivalence result}

We now separate an algebraic compatibility result from the substantive
estimating-equation connection.  The first result reproduces
CB\textsuperscript{2} after its weighted group summaries have already been
computed; because a saturated two-cell GLS can represent any pair of
independent group summaries, this result should not be interpreted as proof
that the implemented \texttt{bbreg()} estimator strictly nests the original
estimator.  For the original test, let
\[
 \widehat p_A=\sum_{i\in A}a_iY_i,\qquad
 \widehat p_B=\sum_{i\in B}b_iY_i,
 \qquad \sum_{i\in A}a_i=\sum_{i\in B}b_i=1,
\]
where the depth- and dispersion-dependent weights $a_i,b_i$ are those
estimated by CB\textsuperscript{2}.  Let its estimated group variances be
$V_A>0$ and $V_B>0$.  The original guide statistic and reference degrees of
freedom are
\begin{align}
 T_{\mathrm{CB2}} &=
 \frac{\widehat p_B-\widehat p_A}{\sqrt{V_A+V_B}}, \label{eq:legacyt}\\
 \nu_{\mathrm{CB2}} &=
 \frac{(V_A+V_B)^2}
 {V_A^2/(n_A-1)+V_B^2/(n_B-1)}. \label{eq:legacydf}
\end{align}

\begin{lemma}[Legacy CB\textsuperscript{2} representation]
\label{lem:legacy}
Let $X$ have row $(1,0)$ for samples in group $A$ and row $(1,1)$
for samples in group $B$.  Define the working precision matrix
\[
 \Omega=\operatorname{diag}\left(
 \left\{\frac{a_i}{V_A}:i\in A\right\},
 \left\{\frac{b_i}{V_B}:i\in B\right\}\right).
\]
Then the identity-link generalized least-squares estimator satisfies
\[
 \widehat{\boldsymbol\beta}_{\mathrm{GLS}}
 =(X^\top\Omega X)^{-1}X^\top\Omega\mathbf{Y}
 =
 \begin{pmatrix}\widehat p_A\\ \widehat p_B-\widehat p_A\end{pmatrix},
\]
and
\[
 (X^\top\Omega X)^{-1}
 =
 \begin{pmatrix}
 V_A & -V_A\\
 -V_A & V_A+V_B
 \end{pmatrix}.
\]
Consequently, the GLS $t$ statistic for the slope is exactly
$T_{\mathrm{CB2}}$.
\end{lemma}

\begin{proof}
Because the weights within each group sum to one, their groupwise total
working precisions are
\[
 \sum_{i\in A}\frac{a_i}{V_A}=\frac{1}{V_A},
 \qquad
 \sum_{i\in B}\frac{b_i}{V_B}=\frac{1}{V_B}.
\]
Hence
\[
 X^\top\Omega X=
 \begin{pmatrix}
 V_A^{-1}+V_B^{-1} & V_B^{-1}\\
 V_B^{-1} & V_B^{-1}
 \end{pmatrix},
\quad
 X^\top\Omega\mathbf{Y}=
 \begin{pmatrix}
 \widehat p_A/V_A+\widehat p_B/V_B\\
 \widehat p_B/V_B
 \end{pmatrix}.
\]
Direct inversion gives the displayed covariance matrix and multiplication
gives the displayed coefficient vector.  Its slope standard error is
$\sqrt{V_A+V_B}$, proving the result.
\end{proof}

\begin{corollary}[Legacy-compatibility identity]
If the slope statistic in Lemma~\ref{lem:legacy} is referred to
$t_{\nu_{\mathrm{CB2}}}$ using Equation~\eqref{eq:legacydf}, its two-sided
$p$-value is identical to the original CB\textsuperscript{2} $p$-value.
Thus identity-link GLS algebra reproduces the original two-group inference on
its group-summary scale.  This corollary is a compatibility check, not a claim
of finite-sample equality with the default logit \texttt{bbreg()} fit.
\end{corollary}

The default implementation uses the logit link because it respects the
parameter space for arbitrary designs.  It is therefore important to prove a
second, weaker relationship rather than claim finite-sample identity.

\begin{proposition}[First-order equivalence of the logit contrast]
\label{prop:local}
Suppose the two groups are independent, their weighted proportion estimators
are consistent, their groupwise and pooled precision estimators converge to a
common $\kappa$, and under the null or a sequence of local alternatives
$p_A,p_B\rightarrow p\in(0,1)$ with
$\widehat p_B-\widehat p_A=O_p(n^{-1/2})$.  Let
$h(p)=\logit(p)$ and define the delta-method statistic
\[
 T_{\logit}=
 \frac{h(\widehat p_B)-h(\widehat p_A)}
 {\sqrt{h'(\widehat p_A)^2V_A+h'(\widehat p_B)^2V_B}}.
\]
Then
\[
 T_{\logit}-T_{\mathrm{CB2}}\xrightarrow{p}0.
\]
Under a correctly specified common-dispersion beta-binomial model, the
binary-design IRLS Wald statistic has the same first-order expansion.  If
$n_A,n_B\rightarrow\infty$ through increasing numbers of independent
biological libraries, both its residual degrees of freedom and
Equation~\eqref{eq:legacydf} diverge, so their reference distributions also
converge to the same standard normal limit.  Holding $n_A,n_B$ fixed while
only increasing read depths $N_i$ is not this asymptotic regime.
\end{proposition}

\begin{proof}
The mean-value theorem and consistency give
\[
 h(\widehat p_B)-h(\widehat p_A)
 =h'(p)(\widehat p_B-\widehat p_A)+o_p(n^{-1/2}),
\]
where $h'(p)=\{p(1-p)\}^{-1}$.  Continuity of $h'$ gives
\[
 \sqrt{h'(\widehat p_A)^2V_A+h'(\widehat p_B)^2V_B}
 =h'(p)\sqrt{V_A+V_B}\{1+o_p(1)\}.
\]
The common positive factor $h'(p)$ cancels, yielding
$T_{\logit}=T_{\mathrm{CB2}}+o_p(1)$.  It remains to connect this transformed
cell-mean statistic to IRLS.  Write $\rho=1/(\kappa+1)$.  Within a group having
mean $\mu_g$, the beta-binomial quasi-score for its logit cell mean is
proportional to
\[
 \sum_{i\in g}\frac{N_i}{\kappa+N_i}(Y_i-\mu_g).
\]
The original CB\textsuperscript{2} weight is
\[
 \left(\frac{1}{\kappa}+\frac{1}{N_i}\right)^{-1}
 =\frac{\kappa N_i}{\kappa+N_i}.
\]
The factor $\kappa$ disappears on normalization, so for fixed common
dispersion the two methods use identical within-group relative weights.
Under the stated common-dispersion model, the separately estimated legacy
precisions and the IRLS precision converge to the same $\kappa$; the resulting
binary-design cell means therefore have the same first-order expansion.
Finally, both $t$ distributions approach the standard normal as their degrees
of freedom diverge.
\end{proof}

These results establish two deliberately separate claims.  The original
statistic has an \emph{algebraically exact legacy representation} after its
group summaries are computed, while the implemented logit statistic is
\emph{locally asymptotically equivalent} in the unadjusted binary design under
the stated common-dispersion assumptions.  Strict finite-sample nesting of the
implemented estimator is not claimed.  In finite samples,
\texttt{bbreg(\textasciitilde group)} need not equal
\texttt{measure\_sgrna\_stats()}: the effect scale, dispersion pooling,
working weights, and degrees of freedom differ.  What regression adds is an
estimable coefficient for continuous predictors, adjustment variables, and
interactions; what it changes must be assessed by calibration and benchmark
data rather than hidden behind the word ``generalization.''  The executable
example \texttt{examples/cb2\_generalization\_proof.R} verifies
Lemma~\ref{lem:legacy} to machine precision and illustrates the
convergence in Proposition~\ref{prop:local}.  In its deterministic guide,
original CB\textsuperscript{2} and the GLS contrast both give effect
$9.059929\times10^{-4}$, standard error $1.113105\times10^{-4}$,
$t=8.139330$, $\nu=5.647255$, and two-sided
$p=2.512928\times10^{-4}$; the maximum numerical discrepancy is
$1.78\times10^{-15}$.  In the link-scale illustration, shrinking the
two-group proportion difference from $8\times10^{-4}$ to
$2.5\times10^{-5}$ reduces the absolute gap between raw and logit statistics
from $0.0583$ to $0.000151$.
